\documentclass[10pt]{article}
\usepackage{hyperref}
\usepackage{amsmath, amssymb, amsfonts}
\usepackage[margin=1cm]{geometry}
\usepackage{xcolor}
\usepackage{graphicx}
\usepackage{enumitem}
\usepackage{multicol}
\usepackage{titlesec}
\parindent 0pt
\setlength{\columnsep}{18pt}
\setlist{noitemsep, topsep=1pt, leftmargin=1.4em}
\titlespacing*{\section}{0pt}{3pt}{1pt}
\titleformat{\section}{\normalfont\large\bfseries}{\thesection.}{0.5em}{}

\begin{document}
\begin{center}
  {\large\bfseries RF and Microwave Engineering (EX752 / EX716)}\\[0.2em]
  {\normalsize Concise PYQ}
\end{center}
\vspace{-0.4em}
\small
\textbf{Notes:} TU IOE PYQ 2069--2082, 24 papers. Regular = \textbf{bold}, Back = normal; ``Regular / Back'' counts as Regular.
BEX (EX752 / EG785, IV/II) = normal font, BEI (EX716, IV/I) = \texttt{mono}.
Marks in (); unique values only.
Keys: (I)=Importance/Need \quad (+)=Advantage \quad (-)=Disadvantage \quad +Fig=Figure \quad +Eg=Example \quad Vs=Compare/Differentiate.
\vspace{2pt}
\hrule
\vspace{3pt}

\begin{multicols}{2}
\footnotesize

\section*{1. Introduction}
\begin{itemize}
  \item Microwave frequency bands (IEEE) + applications; comm protocols of two / three / five systems \hfill (4)(2+4)(3+3)(2+6)(4+6)(5+5)(3+3+4)
  \item Conventional Vs RF/microwave circuits; lumped Vs distributed; passive components at low Vs microwave bands \hfill (4+2)(3+3)(8)(10)
  \item (+)(-) of microwaves Vs acoustic / seismic / conventional / lower-freq \hfill (5)(6)(2+4)(4+4)
  \item Why S-parameter based analysis is a must in microwave network analysis \hfill (5+3)
\end{itemize}

\section*{2. RF \& M/W Transmission Lines}
\begin{itemize}
  \item Microwave TL Vs conventional; microstrip Vs strip-lines; TL behaviour at low Vs microwave bands \hfill (4)(5)(10)
  \item Admittance / immittance chart; sketch + duality parameters + compare scales \hfill (2)(10)(4)(5)
  \item Smith Chart on $50\Omega$ line with $75{+}j100$: $\Gamma_L$, VSWR, $Z_{in}$@$0.375\lambda$, shortest purely-resistive length \hfill (10)
  \item Single-stub matching ($Z_L{=}80{+}j100$; $\Gamma_L{=}0.5\angle51^\circ$; $120{-}j160$/100$\Omega$; $78.27{+}j60.93$/75$\Omega$) +Smith, +S-matrix \hfill (2+8)(8)(5+3+2)(4+4+2)
  \item Double-stub matching ($3\lambda/8$) +Fig/+S-matrix; various loads \hfill (8)(10)(12)(8+2)(10+2+2)
    \begin{itemize}
      \item $Z_L{=}300{+}j300$, $80{+}j180$@3GHz, $190{+}j110$@10GHz, $110{+}j110$/100$\Omega$
      \item $60{-}j80$/50$\Omega$, $109{+}j120$/75$\Omega$ coax, $\Gamma{=}0.64\angle58^\circ$/100$\Omega$, $\Gamma{=}0.45\angle60^\circ$@10GHz, $\Gamma{=}0.33\angle66^\circ$/75$\Omega$ patch
      \item admittance $0.00813{+}j0.0065\,\mho$/50$\Omega$; short-ckt $75{+}j75$, $0.4{+}j0.85$ \hfill (3+15)(10+2)
      \item (+) of double over single stub; matching techniques + S-matrix of perfect match \hfill (2+8)(4+6)
    \end{itemize}
\end{itemize}

\section*{3. RF \& M/W Network Theory}
\begin{itemize}
  \item S-parameters (I); Vs h-parameters; return \& insertion loss \hfill (4)(4+4)(2)(5+5)
  \item S-parameters --- two-port (derive); three-port network + characteristic parameters +Eg \hfill (2+6)(4+5)(4+4)(4+6)(5+2)
  \item S-matrix analysis: reciprocal / lossless / return loss / reflected power; S-matrix of an Amplitude Modulator \hfill (3+1+1+1+2)(3+7)
  \item Differences between two given S-Matrix sets \hfill (4)
\end{itemize}

\section*{4. RF/Microwave Components \& Devices}
\begin{itemize}
  \item Waveguide --- rect TE/TM field equations; modes + parameters \hfill (7+3)(8+2)(8)(10)
    \begin{itemize}
      \item $\text{TE}_{10}$ Vs $\text{TE}_{20}$ cutoff; prove $\text{TM}_{01}/\text{TM}_{10}$ absent; lowest TM cutoff is $\text{TM}_{11}$; $\text{TE}_{10}$ dominant when $b{>}a$ \hfill (2)(3)
      \item dominant Vs degenerate modes; dominant modes (rect \& circular) + 2.254\,cm WG at 6\,GHz?; $a{=}3b$ dominant mode; rect WG (assumed dims): cutoff, $v_p$, $\lambda_g$, $Z_0$ \hfill (2+2)(1+1+2)(6)
      \item circular waveguide TE/TM field; air-filled rect (cutoff/velocities/$Z_0$) +directional coupler \hfill (10)(7)(6+4)
    \end{itemize}
  \item Tee junctions --- Magic Tee (S-params, +derive); E-plane Vs H-plane; hybrid tee (terminated E/H/co-planar arms) \hfill (3+3)(4)(5)(8)(4+4+2)
    \begin{itemize}
      \item two Magic Tees joined at E-arms (derive S-matrix); combined effect in H-plane; magic tee with both arms shorted \hfill (10)(5)(4)
    \end{itemize}
  \item Passive device from S-matrix (4-port); power splitter; E/H-Tee \& Duplexer S-matrices \hfill (5)(6)(8)(10)
  \item Two radar transmitters needing double power; ASR duplexer at half input power \hfill (6)(8)(10)
  \item Directional couplers; circulators; cavity resonators; microstrips \hfill (6)(5)(4)(5)
  \item Probe- Vs loop-coupling; Gunn diode; MASER \hfill (5)(4)
\end{itemize}

\section*{5. Microwave Generators}
\begin{itemize}
  \item Transit-time effect + two-cavity klystron; bunching (density modulation); bunching + field effect \hfill (2+8)(2)(2+6)(2+2+4)
  \item Klystron oscillator/amplifier (multi-cavity); cross-field effect; velocity modulation + reflex klystron +Fig \hfill (2+7)(2+8)(4)(5)(10)
  \item Magnetron construction/operation; cross-field ($45^\circ/60^\circ/90^\circ$ phase) as power amp; circular beam-type cross-field amplifier \hfill (2+8)(10)(6+2)(8)(3+7)(7)
  \item TWT construction \& operation; conventional-tube limits at microwave; bunching + BWO; slow backward wave structure + LNA \hfill (2+6)(2+8)(5)
\end{itemize}

\section*{6. RF Design Practices}
\begin{itemize}
  \item Filter --- insertion loss method; Butterworth/Chebyshev prototype \hfill (5)(8)(4)(10)
  \item LPF/HPF/BPF via microstrip ($\pi$-section, T-type, series-arm, shunt-arm, double-section, triple-pad) +prototyping \hfill (3+5)(6+4)(10)(7+3)(2+8)(8+2)(5+3)(6+2)(3+8+3)
  \item Amplifier --- stability + max gain (bilateral \& unilateral), +Smith, +flowchart \hfill (5+5)(10)(12)(15)(18)(7+8)(8+4)(4+4+4+4)(6+3+3)
    \begin{itemize}
      \item S-params: $0.894\angle{-60.6}$; $0.55\angle150$; $0.656\angle146.7$; $0.72\angle{-116}$ (GaAs); $0.64\angle{-169}$; $0.55\angle{-160}$ (MESFET); $0.45\angle163$ (GaAs FET); $0.45\angle{-160}$; $0.78\angle{-113}$; $0.38\angle{-158}$ ($Z_s{=}25$, $Z_L{=}40$); $0.6\angle{-140}$; $0.55\angle{-150}$
      \item flowchart (GaAsFET) + $\Gamma_{in}/\Gamma_{out}$ trace; amplifier Vs oscillator (stability factor) \hfill (4)(5+5)(3+8)
    \end{itemize}
  \item Oscillator \& mixer theory \hfill (5)(8)
\end{itemize}

\section*{7. Microwave Antennas \& Propagation}
\begin{itemize}
  \item RF/MW radiation hazards + fields; safety practices \& standards (SARPs, IEEE, EMR) \hfill (8)(5)(6)(3+3)(5+5)
  \item SAR (definition, effect, non-ionizing); how radiation is hazardous; radiation zones of a microwave oven \hfill (6)(10)(4+4)(3+3)
  \item GSM BTS: propagation fields + radiation parameters; 900 Vs 1800 MHz fields; hazardous/safe zone @2600 MHz, 0.5 m antenna \hfill (4+4)(4+2+2)
  \item Reactive Vs near- Vs far-field; antenna as DUT --- measurement tool + working principle; HERP \hfill (5+5)(8)(4)
\end{itemize}

\section*{8. RF/Microwave Measurements}
\begin{itemize}
  \item Power measurement --- VSWR meter + low power / bolometry; VSWR $>$ 10; measuring 7.5 mW; ASR power tool \hfill (2+8)(4+6)(2+6)(10)(8)
  \item Calorimeter --- static / dry / circulating Vs flow; calorimeter wattmeter; static + dynamic circulating \hfill (5)(6)(7)(2+8)(6+2)(10)(6+3)
  \item Bolometer bridge (single-bridge limits, double-channel); bolometry principle; network analyzer (VNA); two-port model \hfill (8)(4)(10)(5)
  \item Microwave Vs low-frequency measurements; spectrum analyzer; device to map BTS radiation zones \hfill (3)(4)(5)(8)
\end{itemize}

\end{multicols}
\end{document}
